% Clase 9 --- V(r) y E(r) del cascarón, uno debajo del otro (§3.3).
% V es CONTINUO pero con un quiebre en r = R; E es DISCONTINUO ahí. Las dos
% cosas son la misma: E = -dV/dr, y la derivada de un quiebre salta.
\begin{center}
\begin{tikzpicture}

\begin{axis}[
  fisfig,
  name=vplot,
  width=11cm, height=4.4cm,
  ylabel={$V(r)$},
  xmin=0, xmax=3.6, ymin=0, ymax=1.35,
  xtick={1}, xticklabels={$R$},
  ytick={1}, yticklabels={$\frac{Q}{4\pi\varepsilon_0 R}$},
  xticklabel style={font=\footnotesize},
]
  \addplot[figblue, line width=1.3pt, domain=0:1, samples=2] {1};
  \addplot[figblue, line width=1.3pt, domain=1:3.55, samples=120] {1/x};
  \draw[figgray, dashed, line width=0.6pt] (axis cs:1,0) -- (axis cs:1,1);
  \node[figlblaux, anchor=west] at (axis cs:1.1,1.12) {continuo, con \textbf{quiebre}};
\end{axis}

\begin{axis}[
  fisfig,
  at={(vplot.below south west)}, anchor=above north west,
  width=11cm, height=4.4cm,
  xlabel={$r$}, ylabel={$E(r)$},
  xmin=0, xmax=3.6, ymin=0, ymax=1.35,
  xtick={1}, xticklabels={$R$},
  ytick={1}, yticklabels={$\frac{Q}{4\pi\varepsilon_0 R^{2}}$},
  xticklabel style={font=\footnotesize},
]
  \addplot[figred, line width=1.3pt, domain=0:1, samples=2] {0};
  \addplot[figred, line width=1.3pt, domain=1:3.55, samples=120] {1/(x*x)};
  \draw[figred, dashed, line width=1.1pt] (axis cs:1,0) -- (axis cs:1,1);
  \node[figlblaux, anchor=west] at (axis cs:1.1,1.12) {\textbf{salta} en $r=R$};
\end{axis}

\end{tikzpicture}\\[0.5em]
{\small\itshape Adentro el cascarón es un \textbf{volumen equipotencial}: $V$
constante, y por eso $E=0$ (la derivada de una constante). Afuera se comporta
como una carga puntual. Lo importante es la relación entre los dos gráficos:
$E = -dV/dr$, así que el \textbf{quiebre} de $V$ en $r=R$ ---donde la pendiente
pasa de $0$ a un valor no nulo--- es exactamente el \textbf{salto} de $E$.
Potencial continuo y campo discontinuo no son dos hechos, sino uno.}
\end{center}
